\documentclass[tikz,border=6pt]{standalone}
\usepackage{ctex}
\usepackage{amsmath}
\usetikzlibrary{calc, arrows.meta}

\begin{document}
\begin{tikzpicture}[scale=1.0, thick]

% 四栏：直线与抛物线 4 种位置
% 抛物线 y^2 = 4x，顶点 (0,0)，焦点 (1,0)
% 各栏宽 4.0cm

% ── 栏 1：相离（直线在抛物线左侧）──
\begin{scope}[xshift=0cm]
  \node[font=\small\bfseries] at (1.6, 4.2) {相离};

  % 抛物线 y^2=4x，画 x in [0,3]
  \draw[black, thick, domain=0:3, samples=60, smooth]
    plot ({(\x)^2/4}, \x);
  \draw[black, thick, domain=-3:0, samples=60, smooth]
    plot ({(\x)^2/4}, \x);

  % 坐标轴
  \draw[->, gray, thin] (-0.3, 0) -- (3.3, 0) node[right, font=\footnotesize] {$x$};
  \draw[->, gray, thin] (0, -3.5) -- (0, 3.5) node[above, font=\footnotesize] {$y$};
  \node[font=\footnotesize, below right] at (0,0) {$O$};

  % 直线 x = -0.5（在抛物线左侧，无交点）
  \draw[blue, thick] (-0.5, -3.0) -- (-0.5, 3.0);
  \node[blue, font=\small] at (-0.5, 3.3) {$l$};

  \node[font=\footnotesize] at (1.5, -3.9) {$\Delta < 0$，无交点};
\end{scope}

% 分隔线
\draw[gray, thin, dashed] (3.7, -4.2) -- (3.7, 4.6);

% ── 栏 2：相切（直线与抛物线切于一点）──
\begin{scope}[xshift=4.2cm]
  \node[font=\small\bfseries] at (1.6, 4.2) {相切};

  % 抛物线 y^2=4x
  \draw[black, thick, domain=0:3, samples=60, smooth]
    plot ({(\x)^2/4}, \x);
  \draw[black, thick, domain=-3:0, samples=60, smooth]
    plot ({(\x)^2/4}, \x);

  % 坐标轴
  \draw[->, gray, thin] (-0.3, 0) -- (3.3, 0) node[right, font=\footnotesize] {$x$};
  \draw[->, gray, thin] (0, -3.5) -- (0, 3.5) node[above, font=\footnotesize] {$y$};
  \node[font=\footnotesize, below right] at (0,0) {$O$};

  % 切线：对 y^2=4x，斜率 k 的切线为 y=kx+1/k（需代入 Δ=0）
  % 取 k=1：切线 y=x+1，切点 (1,2)
  \draw[blue, thick] (-0.5, 0.5) -- (2.5, 3.5);
  \node[blue, font=\small] at (2.7, 3.5) {$l$};

  % 切点 (1,2)
  \fill[red] (1, 2) circle [radius=1.8pt];
  \node[font=\footnotesize, right] at (1.1, 2) {$T(1,2)$};

  \node[font=\footnotesize] at (1.5, -3.9) {$\Delta = 0$，切点 $T$};
\end{scope}

% 分隔线
\draw[gray, thin, dashed] (7.9, -4.2) -- (7.9, 4.6);

% ── 栏 3：相交两点（直线穿过抛物线）──
\begin{scope}[xshift=8.4cm]
  \node[font=\small\bfseries] at (1.5, 4.2) {相交（两点）};

  % 抛物线 y^2=4x
  \draw[black, thick, domain=0:3, samples=60, smooth]
    plot ({(\x)^2/4}, \x);
  \draw[black, thick, domain=-3:0, samples=60, smooth]
    plot ({(\x)^2/4}, \x);

  % 坐标轴
  \draw[->, gray, thin] (-0.3, 0) -- (3.3, 0) node[right, font=\footnotesize] {$x$};
  \draw[->, gray, thin] (0, -3.5) -- (0, 3.5) node[above, font=\footnotesize] {$y$};
  \node[font=\footnotesize, below right] at (0,0) {$O$};

  % 直线 y = x - 1（过 (1,0)，斜率1）
  % 代入 y^2=4x: (x-1)^2=4x => x^2-6x+1=0 => x=(6±sqrt(32))/2=3±2√2
  % x1=3+2√2≈5.83（超出画布），x2=3-2√2≈0.17，y2≈0.17-1=-0.83
  % 换斜率 k=2: 直线 y=2x，交 y^2=4x: 4x^2=4x => x(x-1)=0 => x=0或1
  % 交点 (0,0),(1,2)
  % 取斜率合适的线：y = x + 1，代入 (x+1)^2=4x => x^2-2x+1=0 => (x-1)^2=0 切线！
  % 改：y = 0.5x + 2，代入 (0.5x+2)^2=4x => 0.25x^2+2x+4=4x => 0.25x^2-2x+4=0 => x^2-8x+16=0 => (x-4)^2=0 切线
  % 取 y = x - 0.5：代入 (x-0.5)^2=4x => x^2-1x+0.25=4x => x^2-5x+0.25=0
  % x=(5±sqrt(24))/2 = 2.5±√6 ≈ 2.5±2.449 => x1≈4.95, x2≈0.051
  % y1≈4.45, y2≈-0.449
  % 取直线 y = 2x - 3：(2x-3)^2=4x => 4x^2-12x+9=4x => 4x^2-16x+9=0
  % x=(16±sqrt(256-144))/8=(16±sqrt(112))/8=(16±4√7)/8=2±√7/2
  % √7≈2.646, √7/2≈1.323 => x1≈3.323,y1≈3.646; x2≈0.677,y2≈-1.646

  \draw[blue, thick] (-0.2, -3.4) -- (3.2, 3.4);
  \node[blue, font=\small] at (3.4, 3.4) {$l$};

  % 两交点
  \coordinate (P1) at (3.323, 3.646);
  \coordinate (P2) at (0.677, -1.646);
  \fill[black] (P1) circle [radius=1.8pt];
  \fill[black] (P2) circle [radius=1.8pt];
  \node[font=\footnotesize, right] at (P1) {$A$};
  \node[font=\footnotesize, right] at (P2) {$B$};

  % 弦长标注
  \draw[<->, red!70, thin] ($(P1)+(0.25,0)$) -- ($(P2)+(0.25,0)$);
  \node[red!80, font=\footnotesize, right] at (1.9, 1.0) {$|AB|$};

  \node[font=\footnotesize] at (1.5, -3.9) {$\Delta > 0$，两交点 $A,B$};
\end{scope}

% 分隔线
\draw[gray, thin, dashed] (12.1, -4.2) -- (12.1, 4.6);

% ── 栏 4：平行对称轴（必单交点）──
\begin{scope}[xshift=12.6cm]
  \node[font=\small\bfseries, align=center] at (1.6, 4.2) {平行对称轴};

  % 抛物线 y^2=4x
  \draw[black, thick, domain=0:3, samples=60, smooth]
    plot ({(\x)^2/4}, \x);
  \draw[black, thick, domain=-3:0, samples=60, smooth]
    plot ({(\x)^2/4}, \x);

  % 坐标轴
  \draw[->, gray, thin] (-0.3, 0) -- (3.3, 0) node[right, font=\footnotesize] {$x$};
  \draw[->, gray, thin] (0, -3.5) -- (0, 3.5) node[above, font=\footnotesize] {$y$};
  \node[font=\footnotesize, below right] at (0,0) {$O$};

  % 平行 x 轴（对称轴）的直线 y = 1.5，与 y^2=4x 交于 (9/4, 1.5)
  \draw[blue, thick] (-0.3, 1.5) -- (3.3, 1.5);
  \node[blue, font=\small] at (3.5, 1.5) {$l$};

  % 唯一交点 (9/4, 3/2) = (2.25, 1.5)
  \fill[red] (2.25, 1.5) circle [radius=1.8pt];
  \node[font=\footnotesize, above right] at (2.25, 1.5) {$P$};

  % 提示：对称轴方向（虚线）
  \draw[gray, dashed, thin] (-0.3, 0) -- (3.3, 0);

  \node[font=\footnotesize, align=center] at (1.5, -3.9) {必唯一交点 $P$};
  \node[font=\footnotesize, align=center, text=gray] at (1.5, -4.4) {（类比双曲线渐近线）};
\end{scope}

\end{tikzpicture}
\end{document}
